\chapter{Проблемы нативной разработки iOS-приложений}\label{ch:---ios-}

\section{Общие проблемы}\label{sec:ios}
Проблемы указаны с учетом специфики проекта, много проблем связанных с финансами, данная платформа не предполагает коммерциализации, поэтому эти проблемы не описаны.

\begin{enumerate}
    \item \textbf{App Store и политика Apple}
    \begin{itemize}
        \item Строгая модерация: долгие проверки, возможные отклонения по формальным причинам.
    \end{itemize}


    \item \textbf{CI/CD и сборка}
    \begin{itemize}
        \item Требуется macOS для сборки (Xcode)\@.
        \item Работа с сертификатами, provisioning profiles -- сложная и неочевидная.
    \end{itemize}

    \item \textbf{Зависимость от Xcode}
    \begin{itemize}
        \item Частые обновления Xcode и macOS\@.
        \item Нельзя собрать под новую iOS без актуального инструментария.
    \end{itemize}

    \item \textbf{Тестирование}
    \begin{itemize}
        \item Ограничения симулятора: нет камеры, Bluetooth и т.\,д.
        \item TestFlight требует времени и регистрации для внешних тестеров.
    \end{itemize}

\end{enumerate}

\section{Проблемы в современной России}\label{sec:--ios---}

\begin{enumerate}
    \item \textbf{Невозможность прямой публикации из РФ}
    \begin{itemize}
        \item Apple не принимает карты российских банков.
        \item Регистрация разработчика из РФ невозможна.
        \item Обход: регистрация через третьи страны (Казахстан, Грузия, Армения и т.\,д.).
    \end{itemize}

    \item \textbf{Риски блокировок}
    \begin{itemize}
        \item Приложения российских банков, медиа и телекомов уже удалялись.
        \item Возможна блокировка аккаунта и удаление приложения без предупреждения.
    \end{itemize}

    \item \textbf{Отсутствие монетизации в РФ}
    \begin{itemize}
        \item App Store не работает полноценно в России.
        \item Большая часть аудитории не может платить через стандартные механизмы.
    \end{itemize}

\end{enumerate}